% Part: first-order-logic
% Chapter: sequent-calculus
% Section: soundness-identity

\documentclass[../../../include/open-logic-section]{subfiles}

\begin{document}

\olfileid{fol}{seq}{sid}

\olsection{صحت با \usetoken{S}{identity}}

\begin{prop}
$\Log{LK}$ همراه با دنباله‌های آغازین و قواعد همانی صحیح است.
\end{prop}

\begin{proof}
دنباله‌های آغازین به صورت ${} \Sequent \eq[t][t]$ معتبرند، زیرا
برای هر !!{structure}~$\Struct M$، $\Sat{M}{\eq[t][t]}$. (توجه کنید
که فرض می‌کنیم ترم $t$ بسته است، یعنی هیچ متغیری ندارد؛ پس انتساب‌های
متغیر بی‌تأثیرند.)

فرض کنید آخرین استنتاج در !!a{derivation}، $=$ باشد. آنگاه مقدمه
$\eq[t_1][t_2], \Gamma \Sequent \Delta, !A(t_1)$ و نتیجه
$\eq[t_1][t_2], \Gamma \Sequent \Delta, !A(t_2)$ است.
!!a{structure}~$\Struct M$ را در نظر بگیرید. باید نشان دهیم نتیجه
معتبر است؛ یعنی اگر $\Sat{M}{\eq[t_1][t_2]}$ و
$\Sat{M}{\Gamma}$، آنگاه یا $\Sat{M}{!C}$ برای یک
$!C \in \Delta$، یا $\Sat{M}{!A(t_2)}$.

بنا بر فرض استقرا، مقدمه معتبر است. یعنی اگر
$\Sat{M}{\eq[t_1][t_2]}$ و $\Sat{M}{\Gamma}$، آنگاه یا (الف) برای
یک $!C \in \Delta$، $\Sat{M}{!C}$، یا (ب)
$\Sat{M}{!A(t_1)}$. در حالت (الف) کار تمام است. حالت (ب) را در نظر
بگیرید. فرض کنید $s$ انتساب متغیری باشد که
$s(x) = \Value{t_1}{M}$. بنا بر
\olref[syn][ass]{prop:sentence-sat-true}،
$\Sat{M}{!A(t_1)}[s]$. چون $\varAssign{s}{s}{x}$، بنا بر
\olref[syn][ext]{prop:ext-formulas}،
$\Sat{M}{!A(x)}[s]$. از آنجا که
$\Sat{M}{\eq[t_1][t_2]}$، داریم
$\Value{t_1}{M} = \Value{t_2}{M}$ و در نتیجه
$s(x) = \Value{t_2}{M}$. با کاربرد دوبارهٔ
\olref[syn][ext]{prop:ext-formulas}، همچنین
$\Sat{M}{!A(t_2)}[s]$. بنا بر
\olref[syn][ass]{prop:sentence-sat-true}،
$\Sat{M}{!A(t_2)}$.
\end{proof}

\end{document}
